% !TeX program = xelatex
% !TEX encoding = UTF-8
% =====================================================================
%  BAO CAO BAI TAP LON - XU LY NGON NGU TU NHIEN (INT547032)
%  He thong Hoi dap Phap luat Viet Nam dua tren RAG
%  Template ke thua tu bao cao NHOM3_TTHCM (cung sinh vien, cung UEH)
%  Bien dich: XeLaTeX
% =====================================================================
\documentclass[12pt]{article}

% ---------- Trang & le (dung quy cach: A4, tren 2.5, duoi 2.5, trai 3.5, phai 2.0) ----------
\usepackage[a4paper,top=2.5cm,bottom=2.5cm,left=3.5cm,right=2.0cm]{geometry}

% ---------- Font & tieng Viet (XeLaTeX + fontspec) ----------
\usepackage{fontspec}
\setmainfont{Liberation Serif} % font thay the tuong thich kich thuoc voi Times New Roman
\setmonofont{Noto Sans Mono}[Scale=0.85] % ho tro day du dau tieng Viet + ky tu ve khung (box-drawing)

% ---------- Co chu 13pt & dan dong 1.5 ----------
\linespread{1.5}
\makeatletter
\renewcommand\normalsize{%
  \@setfontsize\normalsize{13pt}{15.6pt}%
  \abovedisplayskip 10\p@ \@plus2\p@ \@minus5\p@
  \abovedisplayshortskip \z@ \@plus3\p@
  \belowdisplayshortskip 6\p@ \@plus3\p@ \@minus3\p@
  \belowdisplayskip \abovedisplayskip
  \let\@listi\@listI}
\makeatother
\normalsize

% ---------- Cac goi tien ich chuan, an toan ----------
\usepackage{graphicx}
\setlength{\parskip}{0pt}
\usepackage{indentfirst}
\setlength{\parindent}{1cm}
\usepackage{fancyhdr}
\usepackage[colorlinks=true,linkcolor=black,urlcolor=blue,citecolor=black,linktoc=all]{hyperref}
\usepackage{xurl}
\usepackage{tabularx}
\usepackage{array}
\usepackage{booktabs}
\usepackage{enumitem}
\renewcommand{\labelitemi}{-}
\usepackage{titlesec}
\usepackage{fancyvrb}
\fvset{fontsize=\small,frame=single,framesep=2.5mm,xleftmargin=0pt,xrightmargin=0pt,framerule=0.4pt}

% ---------- Dinh dang so trang: chinh giua ----------
\pagestyle{fancy}
\fancyhf{}
\fancyfoot[C]{\thepage}
\renewcommand{\headrulewidth}{0pt}

% ---------- Muc luc: doi ten thanh MUC LUC ----------
\renewcommand{\contentsname}{MỤC LỤC}

% =====================================================================
%  DINH NGHIA CAC CAP TIEU DE THU CONG
% =====================================================================
\newcounter{seca}
\newcounter{secb}[seca]
\newcounter{secc}[secb]
\newcounter{secd}[secc]

\newcommand{\phanlon}[1]{%
  \clearpage
  \refstepcounter{seca}
  \begin{center}
    {\bfseries\fontsize{15pt}{18pt}\selectfont \Roman{seca}. #1}
  \end{center}
  \addcontentsline{toc}{section}{\Roman{seca}. #1}}

\newcommand{\mucon}[1]{%
  \refstepcounter{secb}%
  \addvspace{16pt plus 2pt minus 2pt}\noindent
  {\bfseries\fontsize{13.5pt}{16pt}\selectfont \arabic{secb}. #1}\par
  \nopagebreak\addvspace{6pt plus 1pt minus 1pt}\nopagebreak
  \addcontentsline{toc}{subsection}{\protect\numberline{\arabic{secb}.}#1}}

\newcommand{\muccon}[1]{%
  \refstepcounter{secc}%
  \addvspace{10pt plus 2pt minus 2pt}\noindent
  {\bfseries #1}\par
  \nopagebreak\addvspace{5pt plus 1pt minus 1pt}\nopagebreak
  \addcontentsline{toc}{subsubsection}{#1}}

\newcommand{\mucnho}[1]{%
  \refstepcounter{secd}%
  \addvspace{8pt plus 1pt minus 1pt}\noindent
  {\bfseries #1}\par\nopagebreak\addvspace{4pt plus 1pt minus 1pt}\nopagebreak}

\newcommand{\refitem}[1]{%
  \par\noindent\hangindent=1.25cm\hangafter=1 #1\par\vspace{6pt}}

% Hinh minh hoa (khong dung float, chen truc tiep tai vi tri goi lenh)
\newcommand{\figblock}[3]{%
  \begin{center}
    \includegraphics[width=#2\linewidth]{#1}\\[4pt]
    {\small\itshape #3}
  \end{center}
  \vspace{4pt}}

% Ghi chu / nhan xet (khung don gian, khong mau)
\newenvironment{ghichu}{%
  \begin{quote}\itshape}
  {\end{quote}}

\begin{document}

% =====================================================================
%  TRANG BIA
% =====================================================================
\begin{titlepage}
\thispagestyle{empty}
\begin{center}
\includegraphics[width=3.0cm]{ueh_logo.png}\\[4pt]
{\fontsize{14pt}{17pt}\selectfont ĐẠI HỌC KINH TẾ THÀNH PHỐ HỒ CHÍ MINH (UEH)}\\[1pt]
{\fontsize{11pt}{14pt}\selectfont TRƯỜNG CÔNG NGHỆ VÀ THIẾT KẾ}\\[2pt]
\rule{0.55\linewidth}{0.6pt}

\vspace{0.6cm}

{\fontsize{13pt}{16pt}\selectfont Môn học: \textbf{XỬ LÝ NGÔN NGỮ TỰ NHIÊN}}\\[2pt]
{\fontsize{12.5pt}{15pt}\selectfont Mã học phần: \textbf{INT547032}}\\[2pt]
{\fontsize{12.5pt}{15pt}\selectfont Lớp học phần: \textbf{26C1INT54703201}}

\vspace{0.7cm}

{\fontsize{12.5pt}{15pt}\selectfont\bfseries BÁO CÁO BÀI TẬP LỚN}\\[10pt]
{\fontsize{16pt}{19pt}\selectfont\bfseries
HỆ THỐNG HỎI ĐÁP PHÁP LUẬT VIỆT NAM\\
DỰA TRÊN RAG\\[4pt]}
{\fontsize{11pt}{14pt}\selectfont\itshape
Vietnamese Legal Question Answering System based on\\
Retrieval-Augmented Generation}

\vspace{0.7cm}

\begin{tabularx}{0.78\linewidth}{@{}lX@{}}
\textbf{Giảng viên hướng dẫn:} & ThS. Nguyễn Khắc Toàn \\[4pt]
\textbf{Sinh viên thực hiện:} & Hồ Minh Hiếu \\[4pt]
\textbf{Mã số sinh viên:} & 31241022078 \\
\end{tabularx}

\vspace{1.0cm}

{\fontsize{13pt}{16pt}\selectfont TP. Hồ Chí Minh, ngày 26 tháng 8 năm 2026}

\end{center}
\end{titlepage}

\clearpage
\pagenumbering{arabic}
\setcounter{page}{1}
\tableofcontents

% =====================================================================
%  I. GIOI THIEU DE TAI
% =====================================================================
\phanlon{GIỚI THIỆU ĐỀ TÀI}

\mucon{Đặt vấn đề}

Hệ thống pháp luật Việt Nam bao gồm hàng nghìn văn bản quy phạm pháp luật — Hiến pháp, Bộ luật, Luật, Nghị định, Thông tư — với tổng khối lượng lên đến hàng triệu từ. Việc tra cứu, tìm hiểu các quy định pháp luật đòi hỏi nhiều thời gian và kiến thức chuyên sâu, gây khó khăn cho người dân, sinh viên và ngay cả các chuyên gia pháp lý.

Với sự phát triển của Xử lý Ngôn ngữ Tự nhiên (NLP) và mô hình ngôn ngữ lớn (LLM), việc xây dựng một hệ thống hỏi đáp pháp luật tự động trở nên khả thi. Người dùng có thể đặt câu hỏi bằng ngôn ngữ tự nhiên và nhận câu trả lời chính xác kèm trích dẫn điều luật cụ thể.

\mucon{Mục tiêu đề tài}

Đề tài xây dựng hệ thống Hỏi đáp Pháp luật Việt Nam dựa trên RAG (Retrieval-Augmented Generation) với các mục tiêu:

\begin{itemize}[leftmargin=1.1cm]
\item Thu thập và xử lý toàn bộ văn bản pháp luật Việt Nam đang có hiệu lực;
\item Xây dựng pipeline RAG end-to-end: chunking $\rightarrow$ embedding $\rightarrow$ retrieval $\rightarrow$ generation;
\item Đánh giá hiệu quả và so sánh các mô hình embedding;
\item Triển khai giao diện web tương tác cho người dùng cuối.
\end{itemize}

\mucon{Phạm vi đề tài}

\begin{itemize}[leftmargin=1.1cm]
\item \textbf{Dữ liệu:} Hiến pháp, Bộ luật, Luật do Quốc hội ban hành (1945--2026);
\item \textbf{Ngôn ngữ:} Tiếng Việt;
\item \textbf{Bài toán:} Question Answering dựa trên tài liệu (Open-domain QA).
\end{itemize}

% =====================================================================
%  II. CO SO LY THUYET
% =====================================================================
\phanlon{CƠ SỞ LÝ THUYẾT}

\mucon{Retrieval-Augmented Generation (RAG)}

RAG (Lewis và cộng sự, 2020) là kiến trúc kết hợp hai bước:

\begin{itemize}[leftmargin=1.1cm]
\item \textbf{Bước 1 — Retrieval (Tìm kiếm):} câu hỏi của người dùng được chuyển thành vector (embedding), sau đó tìm kiếm các đoạn văn bản liên quan nhất trong cơ sở dữ liệu vector thông qua độ tương đồng cosine.
\item \textbf{Bước 2 — Generation (Sinh câu trả lời):} các đoạn văn bản liên quan được đưa vào LLM cùng với câu hỏi để sinh câu trả lời chính xác, có ngữ cảnh.
\end{itemize}

\begin{ghichu}
Ưu điểm của RAG so với fine-tuning thuần túy: không cần train lại mô hình khi thêm tài liệu mới, câu trả lời có thể trích dẫn nguồn cụ thể, và giảm hiện tượng “hallucination” của LLM.
\end{ghichu}

\mucon{Sentence Embeddings}

Sentence embedding là kỹ thuật ánh xạ câu văn thành vector trong không gian chiều cao, sao cho các câu có ý nghĩa tương tự có vector gần nhau (cosine similarity cao).

\begin{itemize}[leftmargin=1.1cm]
\item \textbf{multilingual-e5-base} (Wang và cộng sự, 2024) là mô hình dựa trên kiến trúc Transformer, được pre-train trên 94 ngôn ngữ với kỹ thuật contrastive learning. Mô hình yêu cầu prefix “query:” cho câu hỏi và “passage:” cho đoạn văn bản để tối ưu hóa hiệu suất retrieval.
\item \textbf{vietnamese-sbert} là mô hình SBERT được fine-tune đặc biệt cho tiếng Việt, sử dụng kiến trúc Siamese network để tối ưu hóa độ tương đồng ngữ nghĩa.
\end{itemize}

\mucon{Vector Database}

ChromaDB là cơ sở dữ liệu vector mã nguồn mở, hỗ trợ lưu trữ và truy vấn embedding với thuật toán HNSW (Hierarchical Navigable Small World), cho phép tìm kiếm gần đúng (ANN) hiệu quả với độ phức tạp $O(\log n)$.

\mucon{Chiến lược Chunking}

Để xử lý tài liệu pháp luật, đề tài sử dụng chiến lược chunking theo cấu trúc văn bản:

\begin{itemize}[leftmargin=1.1cm]
\item \textbf{Ưu tiên:} tách theo ranh giới điều luật (Điều 1, Điều 2,...) để giữ nguyên tính toàn vẹn pháp lý;
\item \textbf{Fallback:} tách theo kích thước cố định (512 ký tự) với overlap 50 ký tự khi không có cấu trúc điều.
\end{itemize}

% =====================================================================
%  III. KIEN TRUC HE THONG
% =====================================================================
\phanlon{KIẾN TRÚC HỆ THỐNG}

% CẦN CẬP NHẬT: hình này vẫn vẽ khối "Gradio UI". Kiến trúc thực tế đã đổi thành
% FastAPI (server.py) + giao diện web tĩnh (web/). Hình được tạo bằng công cụ ngoài,
% không có mã sinh trong repo nên không thể sinh lại tự động.
\figblock{fig_architecture.png}{0.92}{Hình 1. Kiến trúc tổng quan hệ thống RAG pháp luật Việt Nam}

\mucon{Hai giai đoạn chính}

\textbf{(A) Offline Indexing} — chạy một lần:

\begin{itemize}[leftmargin=1.1cm]
\item Load 624 văn bản luật từ HuggingFace (\texttt{undertheseanlp/UTS\_VLC});
\item Tách thành 48.803 chunks theo cấu trúc Điều/Khoản;
\item Embed bằng multilingual-e5-base (768 chiều, normalize cosine);
\item Lưu vào ChromaDB với chỉ số HNSW.
\end{itemize}

\textbf{(B) Online Inference} — mỗi khi người dùng hỏi:

\begin{itemize}[leftmargin=1.1cm]
\item Embed câu hỏi với prefix “query:”;
\item Retrieve top-5 chunks có score cao nhất;
\item Gọi Gemini 3.6 Flash với prompt tiếng Việt + context;
\item Trả kết quả kèm trích dẫn tên luật và retrieval score.
\end{itemize}

% =====================================================================
%  IV. DU LIEU VA TIEN XU LY
% =====================================================================
\phanlon{DỮ LIỆU VÀ TIỀN XỬ LÝ}

\mucon{Dataset — UTS\_VLC (Underthesea Vietnamese Legal Corpus)}

\begin{itemize}[leftmargin=1.1cm]
\item Nguồn: \texttt{undertheseanlp/UTS\_VLC} trên HuggingFace;
\item Nội dung: toàn bộ Hiến pháp, Bộ luật, Luật từ 1945--2026, đã được xác minh đang có hiệu lực từ cơ sở dữ liệu chính thức vbpl.vn;
\item Tổng tài liệu: 624 văn bản (splits: 2026, 2023, 2021);
\item Kích thước trung bình: khoảng 370.000 ký tự/tài liệu.
\end{itemize}

\figblock{fig_dataset.png}{0.85}{Hình 2. Phân bố tài liệu theo split và thống kê chunking (UTS\_VLC)}

\begin{center}
\begin{tabularx}{0.92\linewidth}{@{}lccX@{}}
\toprule
\textbf{Split} & \textbf{Số văn bản} & \textbf{Tỷ lệ} & \textbf{Mô tả} \\
\midrule
2026 & 306 & 49,0\% & Verified từ vbpl.vn, đang có hiệu lực \\
2023 & 208 & 33,3\% & Văn bản ban hành đến 2023 \\
2021 & 110 & 17,6\% & Văn bản ban hành đến 2021 \\
\midrule
\textbf{Tổng} & \textbf{624} & \textbf{100\%} & --- \\
\bottomrule
\end{tabularx}\\[4pt]
{\small\itshape Bảng 1. Thống kê dataset UTS\_VLC}
\end{center}

\mucon{Tiền xử lý và Chunking}

Chiến lược chunking theo điều luật, ưu tiên tách theo ranh giới “Điều” bằng biểu thức chính quy:

\begin{Verbatim}
DIEU_PATTERN = re.compile(r"(?=Điều\s+\d+\.)", re.UNICODE)
\end{Verbatim}

Kết quả:

\begin{itemize}[leftmargin=1.1cm]
\item Tổng chunks: 48.803;
\item Phần lớn chunks: 200--800 ký tự (tương ứng 1 điều luật);
\item Chunks overflow: tách theo kích thước 512 ký tự với overlap 50 ký tự.
\end{itemize}

% =====================================================================
%  V. CAI DAT VA THUC NGHIEM
% =====================================================================
\phanlon{CÀI ĐẶT VÀ THỰC NGHIỆM}

\mucon{Môi trường thực nghiệm}

\begin{center}
\begin{tabularx}{0.92\linewidth}{@{}lX@{}}
\toprule
\textbf{Thành phần} & \textbf{Chi tiết} \\
\midrule
OS & Ubuntu 22.04 LTS \\
Python & 3.10.12 \\
Embedding model & intfloat/multilingual-e5-base (278M params) \\
Vector database & ChromaDB 0.5+ (HNSW, cosine similarity) \\
LLM & Gemini 3.6 Flash (Google AI) \\
Framework & sentence-transformers, datasets, fastapi, uvicorn, scikit-learn, pytest \\
\bottomrule
\end{tabularx}\\[4pt]
{\small\itshape Bảng 2. Môi trường thực nghiệm}
\end{center}

\mucon{Cấu trúc dự án}

\begin{Verbatim}
nlp_main/
├── src/
│   ├── ingestion/      # loader.py, chunker.py
│   ├── embeddings/     # embedder.py (multilingual-e5-base)
│   ├── vectorstore/    # store.py (ChromaDB, batch insert)
│   ├── generation/     # generator.py (Gemini 3.6 Flash)
│   └── pipeline/       # rag.py (end-to-end)
├── evaluation/         # evaluator.py, eval_dataset.py (10 câu hỏi)
├── server.py           # FastAPI: /api/ask, /api/ask/stream (SSE)
├── web/                # giao diện: index.html, style.css, app.js, field.js
├── scripts/            # index_documents.py
└── tests/              # 7 file test, 17 test cases
\end{Verbatim}

\mucon{Quy trình thực nghiệm}

Bước 1 — Index toàn bộ dataset vào ChromaDB (khoảng 10 phút):

\begin{Verbatim}
python3 scripts/index_documents.py
# Output: 624 docs -> 48.803 chunks -> embedded & stored
\end{Verbatim}

Bước 2 — Chạy kiểm thử (pytest):

\begin{center}
\begin{tabularx}{0.92\linewidth}{@{}Xlcc@{}}
\toprule
\textbf{File test} & \textbf{Module} & \textbf{Số test} & \textbf{Kết quả} \\
\midrule
test\_loader.py & ingestion/loader & 3 & PASS \\
test\_chunker.py & ingestion/chunker & 4 & PASS \\
test\_embedder.py & embeddings/embedder & 3 & PASS \\
test\_store.py & vectorstore/store & 2 & PASS \\
test\_generator.py & generation/generator & 2 & PASS \\
test\_pipeline.py & pipeline/rag & 2 & PASS \\
test\_evaluator.py & evaluation/evaluator & 1 & PASS \\
\midrule
\textbf{TỔNG} & --- & \textbf{17} & \textbf{17/17 PASS} \\
\bottomrule
\end{tabularx}\\[4pt]
{\small\itshape Bảng 3. Kết quả kiểm thử từng module (pytest)}
\end{center}

Bước 3 — Chạy máy chủ:

\begin{Verbatim}
python -m uvicorn server:app --port 8000
# Running on http://127.0.0.1:8000
\end{Verbatim}

% =====================================================================
%  VI. KET QUA DANH GIA
% =====================================================================
\phanlon{KẾT QUẢ ĐÁNH GIÁ}

\mucon{Bộ câu hỏi đánh giá}

10 câu hỏi pháp luật đa dạng, bao phủ các lĩnh vực dân sự, hôn nhân, hình sự:

\begin{center}
\begin{tabularx}{0.92\linewidth}{@{}cXX@{}}
\toprule
\textbf{\#} & \textbf{Câu hỏi} & \textbf{Từ khóa mong đợi} \\
\midrule
Q1 & Hợp đồng vô hiệu khi nào? & vi phạm, điều kiện, vô hiệu \\
Q2 & Tuổi kết hôn tối thiểu? & 18, tuổi, kết hôn \\
Q3 & Quyền sở hữu tài sản? & sở hữu, tài sản, quyền \\
Q4 & Thời hiệu khởi kiện dân sự? & thời hiệu, khởi kiện, năm \\
Q5 & Bồi thường thiệt hại ngoài HĐ? & bồi thường, thiệt hại, trách nhiệm \\
Q6 & Quyền ly hôn? & ly hôn, vợ, chồng \\
Q7 & Điều kiện di chúc hợp lệ? & di chúc, hợp lệ, điều kiện \\
Q8 & Phân loại tội phạm? & tội phạm, phân loại, nghiêm trọng \\
Q9 & Quyền/nghĩa vụ công dân? & quyền, nghĩa vụ, công dân \\
Q10 & Hình phạt tù tối đa? & tù, năm, hình phạt \\
\bottomrule
\end{tabularx}\\[4pt]
{\small\itshape Bảng 4. Bộ 10 câu hỏi đánh giá pháp luật}
\end{center}

\mucon{Kết quả (multilingual-e5-base, full dataset)}

\begin{center}
\begin{tabularx}{0.92\linewidth}{@{}lXX@{}}
\toprule
\textbf{Metric} & \textbf{Kết quả} & \textbf{Ghi chú} \\
\midrule
Hit Rate & 100\% (10/10) & Retrieve đúng từ khóa kỳ vọng \\
Avg Retrieval Score & 0,8765 & Cosine similarity (thang 0--1) \\
Tổng chunks indexed & 48.803 & 624 văn bản, 3 splits \\
Top-k retrieval & 5 chunks/câu hỏi & --- \\
\bottomrule
\end{tabularx}\\[4pt]
{\small\itshape Bảng 5. Kết quả đánh giá tổng thể với multilingual-e5-base}
\end{center}

\figblock{fig_eval_per_question.png}{0.85}{Hình 3. Retrieval score theo từng câu hỏi (multilingual-e5-base, full dataset 48k chunks)}

\begin{ghichu}
Hit Rate = 100\% cho thấy với 10/10 câu hỏi, các đoạn luật được retrieve đều chứa từ khóa pháp lý kỳ vọng — hệ thống tìm đúng điều luật liên quan. Avg Retrieval Score = 0,8765 (thang 0--1, cosine similarity) phản ánh mức độ tương đồng ngữ nghĩa rất cao giữa câu hỏi và điều luật được lấy về.
\end{ghichu}

% =====================================================================
%  VII. SO SANH MO HINH EMBEDDING
% =====================================================================
\phanlon{SO SÁNH MÔ HÌNH EMBEDDING}

\mucon{Hai mô hình được so sánh}

\begin{center}
\begin{tabularx}{0.92\linewidth}{@{}lccX@{}}
\toprule
\textbf{Mô hình} & \textbf{Kích thước} & \textbf{Ngôn ngữ} & \textbf{Đặc điểm} \\
\midrule
multilingual-e5-base & 278M & 94 ngôn ngữ & Contrastive learning, prefix query/passage \\
vietnamese-sbert & $\sim$110M & Tiếng Việt & Fine-tuned riêng cho tiếng Việt \\
\bottomrule
\end{tabularx}\\[4pt]
{\small\itshape Bảng 6. Thông tin hai mô hình embedding được so sánh}
\end{center}

\figblock{fig_model_comparison.png}{0.85}{Hình 4. So sánh Hit Rate và Avg Retrieval Score giữa hai mô hình embedding}

\mucon{Kết quả so sánh}

\begin{center}
\begin{tabularx}{0.92\linewidth}{@{}lcccX@{}}
\toprule
\textbf{Mô hình} & \textbf{Hit Rate} & \textbf{Avg Score} & \textbf{Dataset} & \textbf{Đánh giá} \\
\midrule
multilingual-e5-base & 100\% & \textbf{0,8765} & Full 48k chunks & Tốt hơn \\
vietnamese-sbert & 100\% & 0,6207 & 100 docs & Thấp hơn \\
\bottomrule
\end{tabularx}\\[4pt]
{\small\itshape Bảng 7. Kết quả so sánh hai mô hình embedding}
\end{center}

\mucon{Phân tích nguyên nhân multilingual-e5-base vượt trội}

\begin{itemize}[leftmargin=1.1cm]
\item \textbf{Quy mô pre-training lớn:} multilingual-e5-base được pre-train trên 94 ngôn ngữ với hàng tỷ cặp câu, giúp mô hình hiểu ngữ nghĩa tiếng Việt sâu hơn dù không fine-tune riêng.
\item \textbf{Cơ chế prefix query/passage:} phân biệt rõ ràng kiểu embedding của câu hỏi và văn bản, tối ưu hóa khoảng cách vector giữa câu hỏi và đoạn văn bản liên quan.
\item \textbf{Ảnh hưởng quy mô dataset:} so sánh trên full dataset (48k chunks) so với subset 100 docs có thể ảnh hưởng đến kết quả của vietnamese-sbert.
\end{itemize}

\begin{ghichu}
Kết luận: đối với bài toán RAG pháp luật tiếng Việt, multilingual-e5-base là lựa chọn tốt hơn, phù hợp với xu hướng các mô hình đa ngôn ngữ quy mô lớn vượt trội hơn mô hình đơn ngôn ngữ nhỏ hơn khi dữ liệu fine-tune còn hạn chế.
\end{ghichu}

% =====================================================================
%  VIII. DEMO HE THONG
% =====================================================================
\phanlon{DEMO HỆ THỐNG}

\mucon{Giao diện}

Hệ thống cung cấp giao diện web tự viết (FastAPI + HTML/CSS/JS thuần) với các thành phần:

\begin{itemize}[leftmargin=1.1cm]
\item \textbf{Input:} ô nhập câu hỏi tự do bằng tiếng Việt;
\item \textbf{Output 1:} câu trả lời được Gemini 3.6 Flash tổng hợp;
\item \textbf{Output 2:} danh sách nguồn luật được sử dụng;
\item \textbf{Output 3:} các điều luật cụ thể được retrieve kèm retrieval score.
\end{itemize}

\muccon{Phân tách giọng nói bằng kiểu chữ}

Câu do mô hình sinh ra được hiển thị bằng chữ không chân (sans-serif, độ đậm 300); nguyên văn điều luật được trích dẫn hiển thị bằng chữ có chân (serif) đặt trong khung viền riêng. Mục đích của lựa chọn này không phải thẩm mỹ mà là giảm rủi ro người đọc nhầm lời diễn giải của mô hình thành lời của văn bản quy phạm pháp luật --- rủi ro lớn nhất của một hệ thống hỏi đáp pháp luật dựa trên sinh văn bản. Màu sắc mang cùng nghĩa: màu cam dành cho truy vấn và nội dung do máy sinh, màu xanh ngọc dành cho trích dẫn luật đã đối chiếu được về điều khoản gốc.

\mucon{Ví dụ kết quả}

\textbf{Câu hỏi:} “Hợp đồng vô hiệu khi nào?”

\begin{ghichu}
Trả lời: Dựa trên các điều luật được cung cấp, hợp đồng bảo hiểm vô hiệu trong các trường hợp vi phạm quy định pháp luật... Đối với hợp đồng lao động, hợp đồng vô hiệu toàn bộ khi...
\end{ghichu}

Nguồn luật:

\begin{itemize}[leftmargin=1.1cm]
\item Luật Kinh doanh bảo hiểm 2000 (24/2000/QH10);
\item Bộ luật Lao động 2019.
\end{itemize}

Retrieval scores: 0,89 · 0,87 · 0,86 · 0,85 · 0,84

\mucon{Chạy demo}

\begin{Verbatim}
cd ~/Desktop/STUDY/nlp_main
python -m uvicorn server:app --port 8000
# Truy cap: http://127.0.0.1:8000
\end{Verbatim}

% =====================================================================
%  IX. KET LUAN
% =====================================================================
\phanlon{KẾT LUẬN}

\mucon{Kết quả đạt được}

Đề tài đã xây dựng thành công hệ thống Hỏi đáp Pháp luật Việt Nam dựa trên RAG với các kết quả:

\begin{center}
\begin{tabularx}{0.92\linewidth}{@{}lXX@{}}
\toprule
\textbf{Hạng mục} & \textbf{Kết quả} & \textbf{Ghi chú} \\
\midrule
Dữ liệu & 624 văn bản luật, 48.803 chunks & UTS\_VLC, 3 splits \\
Hit Rate & 100\% (10/10) & Bộ câu hỏi pháp luật \\
Avg Retrieval Score & 0,8765 & Cosine similarity \\
Test coverage & 17/17 PASS & Toàn bộ test cases \\
Demo & Giao diện web & http://localhost:8000 \\
Model comparison & 2 mô hình & e5-base tốt hơn vi-sbert \\
\bottomrule
\end{tabularx}\\[4pt]
{\small\itshape Bảng 8. Tóm tắt toàn bộ kết quả dự án}
\end{center}

\mucon{Điểm nổi bật}

\begin{itemize}[leftmargin=1.1cm]
\item Chunking thông minh theo Điều/Khoản giúp giữ nguyên tính toàn vẹn pháp lý;
\item Embedding hiệu quả: Hit Rate 100\%, Avg Score 0,8765 trên 48k chunks;
\item Trích dẫn nguồn: mỗi câu trả lời kèm tên luật, mã văn bản và retrieval score;
\item So sánh khách quan 2 embedding model với số liệu thực tế;
\item Full pipeline: từ thu thập dữ liệu đến demo sản phẩm hoàn chỉnh, có test đầy đủ.
\end{itemize}

\mucon{Hướng phát triển}

\begin{itemize}[leftmargin=1.1cm]
\item \textbf{Re-ranking:} sử dụng cross-encoder để cải thiện độ chính xác retrieval;
\item \textbf{Knowledge Graph:} tích hợp đồ thị tri thức (như \texttt{hienlh/vn-legal-rag}) để hiểu quan hệ giữa các điều luật;
\item \textbf{Multilingual QA:} hỗ trợ câu hỏi tiếng Anh, tự động dịch và trả lời;
\item \textbf{Cập nhật tự động:} pipeline tự động đồng bộ khi có văn bản luật mới từ vbpl.vn.
\end{itemize}

% =====================================================================
%  TAI LIEU THAM KHAO
% =====================================================================
\clearpage
\begin{center}
{\bfseries\fontsize{15pt}{18pt}\selectfont TÀI LIỆU THAM KHẢO}
\end{center}
\addcontentsline{toc}{section}{TÀI LIỆU THAM KHẢO}
\bigskip

\refitem{Lewis, P., et al. (2020). \textit{Retrieval-Augmented Generation for Knowledge-Intensive NLP Tasks}. NeurIPS 2020. arXiv:2005.11401.}

\refitem{Wang, L., et al. (2024). \textit{Multilingual E5 Text Embeddings: A Technical Report}. arXiv:2402.05672.}

\refitem{Reimers, N., \& Gurevych, I. (2019). \textit{Sentence-BERT: Sentence Embeddings using Siamese BERT-Networks}. EMNLP 2019.}

\refitem{Underthesea. (2026). \textit{UTS\_VLC: Vietnamese Legal Corpus}. HuggingFace Datasets: \texttt{undertheseanlp/UTS\_VLC}.}

\refitem{Google DeepMind. (2024). \textit{Gemini: A Family of Highly Capable Multimodal Models}.}

\refitem{Chroma. (2024). \textit{Chroma: The AI-native open-source embedding database}. \url{https://www.trychroma.com}}

\refitem{hienlh. (2024). \textit{Vietnamese Legal RAG System with Knowledge Graph}. GitHub: \texttt{hienlh/vn-legal-rag}.}

\refitem{CTU-LinguTechies. (2023). \textit{VN-Law-Advisor: Ứng dụng hỏi đáp pháp luật}. GitHub: \texttt{CTU-LinguTechies/VN-Law-Advisor}.}

\end{document}
